\documentclass[12pt]{amsart}
\usepackage{amsmath,amsthm,amssymb}
\usepackage{dsfont,bm,mathtools}
\usepackage[usenames,dvipsnames,svgnames,table]{xcolor}
\usepackage[top=2.3cm,bottom=2.3cm,left=2cm,right=2cm]{geometry}
\usepackage{hyperref}
\hypersetup{hidelinks}
\allowdisplaybreaks
\emergencystretch=1.5em
\definecolor{ProposalOrange}{RGB}{160,65,0}
\definecolor{RefereePurple}{RGB}{103,34,147}
\newtheorem{thm}{Theorem}[section]
\newtheorem{lem}[thm]{Lemma}
\theoremstyle{remark}
\newtheorem{rem}[thm]{Remark}
\theoremstyle{definition}
\newtheorem{nota}[thm]{Notation}
\definecolor{ResolutionGreen}{RGB}{0,100,55}
\newcommand{\C}{\mathds{C}}
\newcommand{\Q}{\mathds{Q}}
\newcommand{\rv}{\operatorname{rv}}
\newcommand{\Gal}{\operatorname{Gal}}
\newcommand{\Aut}{\operatorname{Aut}}
\newcommand{\RV}{\mathrm{RV}}
\newcommand{\RES}{\mathrm{RES}}
\newcommand{\VF}{\mathrm{VF}}
\newcommand{\K}{\mathds{k}}
\newcommand{\Z}{\mathds{Z}}
\newcommand{\N}{\mathds{N}}
\newcommand{\vrv}{\operatorname{vrv}}
\newcommand{\GL}{\operatorname{GL}}
\newcommand{\dcl}{\operatorname{dcl}}
\newcommand{\pr}{\operatorname{pr}}
\newcommand{\defn}{\mathrm{def}}
\newcommand{\given}{\,\middle|\,}
\newcommand{\Sbase}{\mathbb S}
\newcommand{\RVlang}{\mathcal L_{\mathrm{RV}}}
\numberwithin{equation}{section}
\pagestyle{plain}

\newcommand{\A}{\mathds{A}}
\newcommand{\Gm}{\mathbb G_m}
\newcommand{\Var}{\mathrm{Var}}
\newcommand{\Kres}{K_0(\RES)}
\newcommand{\Bangres}{{!K_0(\RES)}}
\newcommand{\Kraw}{K_0^{\widehat\mu,\mathrm{raw}}(\Var_{\C})}
\newcommand{\Kmon}{K_0^{\widehat\mu}(\Var_{\C})}
\newcommand{\E}{\mathbb E_b}
\newcommand{\Vol}{\operatorname{Vol}}
\newcommand{\vv}{\operatorname{val}}
\newcommand{\LiftingTheorem}{the lifting theorem of Section~3}

\theoremstyle{plain}
\newtheorem*{thm*}{Theorem}

\newcommand{\R}{\mathds R}
\newcommand{\Sfiber}{\mathbb S_b}
\newcommand{\Rbase}[1]{{!K_0(\RES_{#1})}}
\newcommand{\IsVFVG}{(\VF,\Gamma)\text{-generated}}
\newcommand{\gsk}{K_+}
\theoremstyle{plain}
\newtheorem{test}[thm]{Test}
\begin{document}
\setcounter{section}{5}
\setcounter{subsection}{2}
\setcounter{thm}{19}
{\small\noindent Forey--Yin, \emph{Bounded integral and motivic Milnor fiber}.
Section 5 review: R094--R095.
Approved context is black; proposed manuscript text is orange.}
\medskip

\noindent\emph{Proof of Theorem 5.17 (continued).}
Put $g=[\Gm]$.
The preceding relative Gamma-push calculation gives
\[
 \mathbb E_{\mathrm b}([\bm W])
   =\sum_i e_i g^{\ell_i}\mathbb E_{\mathrm b}([\mathcal V_i])
\]
when the monomial packages are defined over the RV base.
For a varying value tuple $\alpha(\tau)$ it also gives
\[
 \mathbb E_{\mathrm b}([T_{\alpha,\mathcal Z}])
      =g^n\mathbb E_{\mathrm b}([\mathcal Z]).
\]
We must now obtain the complete finite witnesses over that base.

{\color{ProposalOrange}
\medskip
\noindent\emph{Finite witnesses for the RES quotient.}
The following object-level criterion supplies the finite
data needed for compactness.

\begin{test}\label{test}
Let $\mathbb T$ be a parameter base and
$P,Q\in\RES_{\mathbb T}$.
The equality $[P]=[Q]$ in $R(\mathbb T)$ holds if and only if
there are $C\in\RES_{\mathbb T}$, finitely many ordered pairs
$(L_j,R_j)$, and an actual $\mathbb T$-definable bijection
\begin{equation}\label{finite:RES:witness}
 P\sqcup C\sqcup\bigsqcup_j R_j
       \ \xrightarrow{\ \sim\ }\
 Q\sqcup C\sqcup\bigsqcup_j L_j .
\end{equation}
For each $j$, the pair $(L_j,R_j)$ is either
\[
 \bigl(\alpha_j^\sharp\times Z_j,\
              (\Gm)^{n_j}\times Z_j\bigr)
\]
or its reverse, where
$Z_j\in\RES_{\mathbb T}$ and
$\alpha_j\in\Gamma^{n_j}$ is $\mathbb T$-definable.
All unions and copies are tagged disjoint unions.
\end{test}

To prove the criterion, an element of the ideal defining
$R(\mathbb T)$ is a finite sum of its generators times
elements of $K_0(\RES_{\mathbb T})$.
Write each coefficient as a difference of two positive
object classes. Absorb those objects into $Z_j$, reversing
a pair for a negative coefficient.
The one-coordinate generators already suffice; tuple
generators give the same ideal by expanding a difference
of products.
Thus $[P]-[Q]=\sum_j([L_j]-[R_j])$ in the group completion.
Equality there means that adding one common positive class
$[C]$ makes the two positive sums equal.
Positive scissor equalities assemble into an actual
definable bijection, giving
(\ref{finite:RES:witness}).
Conversely that bijection and the defining ideal relations
imply $[P]=[Q]$ in $R(\mathbb T)$.
This gives finite object witnesses directly.

\medskip
\noindent\emph{Choose the full packages before preparing the base.}
Return temporarily to $b\in B$ in the original parameter
coordinates, and put $\Sfiber=\Sbase\langle b\rangle$.
Apply Lemma 5.12 to the two RV representatives over
$\Sfiber$, with their zero strata and finite tags retained.
The standing nontriviality of the ground Gamma group
also holds over $\Sfiber$.
Only RV decomposition and the RV retraction are used here;
no fiber VF-integration hypothesis is added.

For one such finite package the charts have matrices $M_i$,
free arities $\ell_i$, residue factors $V_{i,b}$,
Gamma supports $I_{i,b}$ and integers
$e_i=\chi_b(I_{i,b})$, and similarly for the primed family.
Set
\[
 C_{i,b}=V_{i,b}\times(\Gm)^{\ell_i},
 \qquad C'_{i',b}=V'_{i',b}\times(\Gm)^{\ell'_{i'}}.
\]
Writing $a^+=\max(a,0)$ and $a^-=\max(-a,0)$, form the
positive RES objects
\[
 P_b=\bigsqcup_i e_i^+ C_{i,b}
       \ \sqcup\ \bigsqcup_{i'}(e'_{i'})^- C'_{i',b},
\]
\[
 Q_b=\bigsqcup_i e_i^- C_{i,b}
       \ \sqcup\ \bigsqcup_{i'}(e'_{i'})^+ C'_{i',b}.
\]
Here $mC$ means $m$ tagged copies of $C$.
The monomial evaluation and the hypothesis imply
$[P_b]=[Q_b]$ in $R(\Sfiber)$.
Choose the full finite witness of Test~\ref{test} over
$\Sfiber$: the stabilizer $C_b$, the factors $Z_{j,b}$,
the value tuples $\alpha_j(b)$, the orientations of the
pairs, and the graph of the displayed bijection.

Include in the package the representative and label
formulas, their finite-fiber bounds, the monomial source
pieces and coordinate maps, the residue and Gamma factors,
and finite bounds for every RES value support.
Include formulas defining all Gamma-valued parameter
functions and the complete stabilization bijection graph.

The Euler integers also have finite first-order certificates.
After DOAG quantifier elimination, the indicator of each
$I_{i,b}$ is a finite integer combination of indicators
of closed rational polyhedra with $b$-definable Gamma
parameters. Its bounded Euler characteristic is the
same combination of their nonemptiness indicators,
since a nonempty closed polyhedron has bounded Euler
characteristic $1$.
Record those parameter formulas, the indicator identity
and the nonemptiness flags.
Validity of a fixed such certificate is first order.

Validity of the entire fixed finite package is now a
first-order condition on $b$: all arities and bounds are
fixed, and the required set identities, maps and actual
bijections are given by formulas.
These conditions cover every ambient $b\in B$.
Saturation and compactness therefore give a finite
definable partition of $B$ with one package on each piece.
In particular, the matrices, copy numbers, Euler integers,
arity lists and pair orientations are fixed on a piece.

\medskip
\noindent\emph{Prepare the witness formulas themselves.}
Repeat the earlier mixed-base preparation with this larger
finite package.
First eliminate valued-field quantifiers in the formulas
for all RV/Gamma sets and maps in it.
Their only free VF variables are the VF coordinates of
the base parameter $b$. Apply the preparation lemma to
every VF polynomial occurring in these formulas, including
those for the stabilization bijection, the Gamma-valued
functions and the Euler certificates.
Carry every original and newly introduced RV coordinate.

As before, write the resulting base as
\[
 B^\flat=\{(c,t):\tau=(\rv(c),t)\in S\}.
\]
Every formula in the RV/Gamma witness package now depends
on $(c,t)$ only through $\tau$.
It therefore defines the representative sets, their label
maps, the monomial data, all auxiliary RES objects,
the Gamma functions $\alpha_j(\tau)$ and the actual
stabilization bijection over $\Sbase\langle\tau\rangle$.
The finite support bounds and the Euler certificates remain
valid there: their defining first-order assertions involve
only $\tau$ after preparation.

This constructs the equality witness over the RV parameter
base. It does not deduce that equality from injectivity of
a base-extension map.
The lift bijections $H_b,H'_b$ are kept as transported
bijections of total mixed VF families; they need not
descend to RV maps.
Using them gives the total RV objects
$\bm W,\bm W'$ exactly as in the earlier mixed-coordinate
construction.
The Gamma functions and support sets may vary with
$\tau$; their constancy is not required.

\medskip
\noindent\emph{Apply the finite witnesses to the total families.}
For every object family $F_\tau$ use the tagged total
\[
 \mathcal F=\{(\tau,u):\tau\in S,\ u\in F_\tau\}.
\]
Give each total a finite-to-one RV label, for example its
full coordinate tuple when necessary.
These totals are RV objects; no assertion that they are
ground RES objects is needed.
The uniform bijection (\ref{finite:RES:witness}) yields
an actual definable bijection
\[
 \mathcal P\sqcup\mathcal C\sqcup\bigsqcup_j\mathcal R_j
       \ \cong\
 \mathcal Q\sqcup\mathcal C\sqcup\bigsqcup_j\mathcal L_j .
\]
For each $j$, the relative torsor identity gives
\[
 \mathbb E_{\mathrm b}([\mathcal L_j])
       =\mathbb E_{\mathrm b}([\mathcal R_j]),
\]
with either orientation, even for varying
$\alpha_j(\tau)$.
Apply the ungraded bounded retraction to the total
bijection and cancel the common summands additively:
\[
 \mathbb E_{\mathrm b}([\mathcal P])
       =\mathbb E_{\mathrm b}([\mathcal Q]).
\]
The definitions of $P_\tau,Q_\tau$ turn their difference
into
\[
 \sum_i e_i g^{\ell_i}\mathbb E_{\mathrm b}([\mathcal V_i])
   -\sum_{i'}e'_{i'}g^{\ell'_{i'}}
                   \mathbb E_{\mathrm b}([\mathcal V'_{i'}]).
\]
By the preceding family-coefficient formula this is
$\mathbb E_{\mathrm b}([\bm W])
 -\mathbb E_{\mathrm b}([\bm W'])$.
Consequently the two global bounded values agree.
Since $\bm W,\bm W'$ represent the original total
VF objects under the lifting comparison, this proves
the conclusion of Theorem 5.17.
Sum over the finitely many base pieces.
All cancellations above take place in the additive
group of $R(\Sbase)$; no torus class is inverted.
}

\medskip
{\normalfont\small\color{RefereePurple}
\noindent\textbf{Referee R094 [Typo; revised] --
Height convention.}
In a construction-graph presentation, height is the
maximum level number, or the number of levels minus one.
The finite object-witness criterion above replaces that
presentation and its height induction.\par
\medskip
\noindent\textbf{Referee R095 [Gap; revised] --
Descent from VF parameters to RV parameters.}
Prepare every formula in the full finite witness package
over the original $b$, including the actual bijections
and Gamma data. Preparing only the object formulas does
not descend the equality witnesses. The proposed text
carries out that larger preparation and forms tagged
total families.\par}
\medskip
{\normalfont\small\color{ProposalOrange}
\noindent\textbf{Proposed disposition of R094--R095.}
Replace the old congruence test and construction-graph
block through the end of the Fubini proof by the orange
text. The following corollary will be reviewed separately.}
\end{document}
